%%% LaTeX Template: Article/Thesis/etc. with colored headings and special fonts
%%%
%%% Source: http://www.howtotex.com/
% vim: set spell spelllang=es syntax=tex :

\documentclass[12pt]{article}
\usepackage{styles/apuntes-estilo}
\usepackage{styles/egyptian}
\usepackage{fancyhdr,lastpage}
\usepackage{hyperref}
\usepackage[inline]{enumitem}
\usepackage{lipsum}

\def\maketitle{

\makeatletter{
    \color{bl} \centering \huge \sc
    \textbf{
        Trabajo práctico obligatorio N° 1\\
        \large \vspace*{-8pt} \color{black}
        ~
        \vspace*{8pt}
    }\\
    \small Fecha de entrega: 22 de Mayo
    \par
}

\makeatother

\makeatletter
% vim: set spell spelllang=es syntax=tex :
 {\centering \small 
    Introducción a la computación\\
    Departamento de Ingeniería de Computadoras \\
    Facultad de Informática - Universidad Nacional del Comahue \\
    \vspace{20pt} }
\makeatother

\vspace{-2.5cm}
\mbox{\hspace{-1cm}\includegraphics[width=3cm,height=3cm]{logos/uncoma.pdf}\hspace{12cm}
    \includegraphics[width=3cm,height=3cm]{logos/fai.pdf}}



}
%\def\maketitle{
%
%\makeatletter{
%    \color{blue} \centering \huge \sc
%    \textbf{
%        Trabajo práctico obligatorio N° 1\\
%        \large \vspace*{-8pt} \color{black}
%        \vspace*{8pt}
%    }\\
%    \small Fecha de entrega: 22 de Mayo
%    \par
%}
%
%\makeatother
%
%\makeatletter
%% vim: set spell spelllang=es syntax=tex :
 {\centering \small 
    Introducción a la computación\\
    Departamento de Ingeniería de Computadoras \\
    Facultad de Informática - Universidad Nacional del Comahue \\
    \vspace{20pt} }
\makeatother

\vspace{-2.5cm}
\mbox{\hspace{-1cm}\includegraphics[width=3cm,height=3cm]{logos/uncoma.pdf}\hspace{12cm}
    \includegraphics[width=3cm,height=3cm]{logos/fai.pdf}}


%
%}

% Custom headers and footers
\fancyhf{} % clear all header and footer fields
\fancypagestyle{plain}{\fancyhf{}}
\pagestyle{fancy}
\lhead{\footnotesize Trabajo práctico obligatorio}
\rhead{\footnotesize \thepage\ }

\def\ti#1#2{\texttt{#1} & #2 \\ }

\begin{document}

\thispagestyle{empty}
\maketitle
\setlength{\parindent}{1pt}


\textbf{Condiciones de la entrega:} \lipsum[1-1]

\begin{enumerate}

    \item Sean $A=1100\,1101$ y $B=0011\,0011$ números enteros con signo
        representados en complemento a dos. Realizar las siguientes
        operaciones:

        \begin{enumerate}

            \item $A+B$

            \item $A-B$

        \end{enumerate}

    \item ¿Cual es la representación en \emph{punto fijo} y \emph{complemento
        a 2} del número $17,75$ si se utilizan 6 bits para la parte entera y 2
        para la parte fraccionaria?

    \item ¿Que número esta representado en los bytes $0x41bb0000$ si la
        representación es \emph{Punto Flotante IEEE-754 de precisión simple
        (32 bits)}?

    \item ¿Que mensaje esta codificado en la siguiente secuencia de bytes
        expresados en hexadecimal: \textbf{41 53 43 49 49 20 65 73  20 69 6e
        74 65 72 65 73 61 6e 74 65 3f}?

\textbf{Tabla \emph{ASCII}:}
\begin{verbatim}
Dec Hex    Dec Hex    Dec Hex  Dec Hex  Dec Hex  Dec Hex   Dec Hex   Dec Hex  
  0 00 NUL  16 10 DLE  32 20    48 30 0  64 40 @  80 50 P   96 60 `  112 70 p
  1 01 SOH  17 11 DC1  33 21 !  49 31 1  65 41 A  81 51 Q   97 61 a  113 71 q
  2 02 STX  18 12 DC2  34 22 "  50 32 2  66 42 B  82 52 R   98 62 b  114 72 r
  3 03 ETX  19 13 DC3  35 23 #  51 33 3  67 43 C  83 53 S   99 63 c  115 73 s
  4 04 EOT  20 14 DC4  36 24 $  52 34 4  68 44 D  84 54 T  100 64 d  116 74 t
  5 05 ENQ  21 15 NAK  37 25 %  53 35 5  69 45 E  85 55 U  101 65 e  117 75 u
  6 06 ACK  22 16 SYN  38 26 &  54 36 6  70 46 F  86 56 V  102 66 f  118 76 v
  7 07 BEL  23 17 ETB  39 27 '  55 37 7  71 47 G  87 57 W  103 67 g  119 77 w
  8 08 BS   24 18 CAN  40 28 (  56 38 8  72 48 H  88 58 X  104 68 h  120 78 x
  9 09 HT   25 19 EM   41 29 )  57 39 9  73 49 I  89 59 Y  105 69 i  121 79 y
 10 0A LF   26 1A SUB  42 2A *  58 3A :  74 4A J  90 5A Z  106 6A j  122 7A z
 11 0B VT   27 1B ESC  43 2B +  59 3B ;  75 4B K  91 5B [  107 6B k  123 7B {
 12 0C FF   28 1C FS   44 2C ,  60 3C <  76 4C L  92 5C \  108 6C l  124 7C |
 13 0D CR   29 1D GS   45 2D -  61 3D =  77 4D M  93 5D ]  109 6D m  125 7D }
 14 0E SO   30 1E RS   46 2E .  62 3E >  78 4E N  94 5E ^  110 6E n  126 7E ~
 15 0F SI   31 1F US   47 2F /  63 3F ?  79 4F O  95 5F _  111 6F o  127 7F DEL
\end{verbatim}

    \item Dada la siguiente imagen y el formato de representación estudiado en
        la materia:

    \includegraphics[width=0.50\textwidth]{img/img_cara.pdf}

    \begin{enumerate}

        \item Representar la imagen utilizando el esquema sin compresión.

        \item Representar la imagen utilizando el esquema de compresión
            \emph{RLE}, utilizando 4 bits para la cantidad.

    \end{enumerate}

\end{enumerate}

\end{document}
